{\footnotesize
\begin{tabularx}{\textwidth}{>{\hsize=0.514\hsize\sloppy\arraybackslash}X >{\hsize=1.486\hsize\sloppy\arraybackslash}X}
\arrayrulecolor{GithubTableBorder}\toprule
\rowcolor{BrandTableHeader}
\textbf{Elemento} & \textbf{Regra} \\
\midrule
\arrayrulecolor{BrandGrayRule}
\texttt{MAJOR} (\texttt{v1.0.0}, \texttt{v2.0.0}, \texttt{v3.0.0}) & Transição entre macroentregas ou mudança estrutural de arquitetura \\
\midrule
\texttt{MINOR} (\texttt{v1.1.0}) & Funcionalidade retrocompatível dentro de uma macroentrega \\
\midrule
\texttt{PATCH} (\texttt{v1.0.1}) & Correção de defeito ou ajuste sem quebra de contrato \\
\midrule
Conventional Commits & \texttt{\textless{}tipo\textgreater{}(\textless{}escopo\textgreater{}):\ \textless{}descrição\textgreater{}}, com tipos \texttt{feat}, \texttt{fix}, \texttt{refactor}, \texttt{test}, \texttt{docs}, \texttt{chore} e \texttt{ci}, para gerar o \emph{changelog} e rastrear as issues \\
\arrayrulecolor{GithubTableBorder}\bottomrule
\end{tabularx}%
}
